\documentclass[11pt,a4paper]{article}
\usepackage[utf8]{inputenc}
\usepackage[margin=1in]{geometry}
\usepackage{amsmath,amssymb}
\usepackage{booktabs}
\usepackage{xcolor}
\usepackage{hyperref}
\usepackage{natbib}

\title{\textbf{Research Scope \& Problematization Memorandum: The Hypocrisy of Parity (Draft)}\\[0.3ex]
\large Discursive Patterns of Victimhood Signaling and Self-Serving Attribution in Exploitative Relationships}
\author{\textbf{Peerapat Asoktummarungsr}}
\date{October 29, 2026 (Draft)}

\begin{document}
\maketitle

\begin{abstract}
This memorandum establishes the formal scope, theoretical problematization, and research gap classification for the manuscript \textit{"The Hypocrisy of Parity: Discursive Patterns of Victimhood Signaling and Self-Serving Attribution in Exploitative Relationships"}. Grounded in Swales' CARS framework and Alvesson \& Sandberg's (2011) problematization methodology, this document maps three intersecting scholarly conversations, classifies the core theoretical paradox within real-world interpersonal and dyadic relationships, articulates three FINER research questions, and conducts a pre-mortem audit to eliminate fatal methodological vulnerabilities prior to full empirical deployment.
\end{abstract}

\section{Scholarly Conversation Mapping}
This study intervenes at the nexus of four intersecting scholarly conversations spanning Discursive Psychology, Interpersonal Communication, Social Psychology, and Computational Social Science:
\begin{enumerate}
    \item \textbf{Conversation 1: Moral Hypocrisy \& Self-Serving Attribution in Relationships.} Foundational works establish that human moral action is frequently guided by the desire to appear fair and principled without incurring the behavioral costs of sacrifice \citep{batson1997unmasking,batson1999moral}. In close relationships, this is reinforced by self-serving attribution and moral licensing, wherein claiming egalitarian values creates cognitive credentials that insulate actors from guilt during subsequent selfish conduct \citep{monin2001moral,merritt2010moral,krettenauer2022when}.
    \item \textbf{Conversation 2: Discursive Psychology, Blame Management, \& Interpersonal Accounts.} Research in discursive psychology and interpersonal communication demonstrates that language in dispute situations functions as a performative social action \citep{edwards1992discursive,billig1996arguing}. Actors systematically deploy conversational accounts, justifications, and excuses \citep{scott1968accounts} to manage accountability, discursively construct victimhood, and deflect responsibility for relationship friction.
    \item \textbf{Conversation 3: Interpersonal Exploitation, Dark Personalities, \& Weaponized Signaling.} Recent research in personality psychology demonstrates that subclinical Dark Triad traits actively deploy virtuous victimhood and dominance-based moral grandstanding to secure relational control, emotional compliance, and unearned partner accommodations \citep{ok2021signaling,bates2024virtuous,cerulli2026dark,savejnarong2022interpersonal}.
    \item \textbf{Conversation 4: Computational Social Science \& Large Language Models as Cultural Informants.} Advances in NLP establish that LLMs can accurately capture pragmatic context, decode subtle interpersonal gaslighting, and identify latent semantic tensions across massive unstructured relationship narratives, bypassing self-report social desirability bias \citep{ziems2024large,argyle2023out,grossmann2023ai,grossmann2024generative}.
\end{enumerate}

\section{Problematization: Beyond Gap-Spotting}
Traditional research on relationship equity and domestic labor relies on simple \textit{gap-spotting} (e.g., investigating whether gender roles are more equal in a specific demographic). Following \citet{alvesson2011generating}, this study pursues \textbf{problematization} by challenging the foundational assumption shared by contemporary relationship science:
\begin{quote}
\textit{Conventional Assumption:} Vocal demands for 50/50 equality, emotional labor sharing, and relationship fairness reflect prosocial dyadic commitment aimed at reducing relational friction.
\end{quote}
\textbf{Problematized Axiom (First Principles):} In interpersonal relationships and the treatment of others, egalitarian vocabulary is structurally optimized to serve as an offensive instrument of interpersonal leverage. Far from promoting equity, high-frequency parity demands combined with victimhood signaling actively lower the perceived moral cost of exploiting others by establishing moral credentials and pre-empting accountability.

\section{FINER Research Questions}
Applying the FINER criteria (\textit{Feasible, Interesting, Novel, Ethical, Relevant}) and safeguarding against Type III errors \citep{alvesson2011generating}, the inquiry is anchored by three core questions:
\begin{itemize}
    \item \textbf{RQ1 (Parity Demands vs. Relational Extraction):} To what extent does high-frequency deployment of Egalitarian Vocabulary \& Parity Demands (\textit{EVP}) correlate with Exploitative Interpersonal Relationship behaviors (\textit{EIR}) in observed relationship dispute disclosures, and is this association mediated by Self-Serving Attribution \& Licensing (\textit{SSA})?
    \item \textbf{RQ2 (Victimhood as an Interpersonal Lever):} How does Virtuous Victimhood Signaling (\textit{VVS}) moderate the relationship between parity rhetoric and self-serving justifications, enabling actors to extract concessions while neutralizing resistance in their treatment of others?
    \item \textbf{RQ3 (AI Cultural Informant Decoding):} How effectively do multi-model Large Language Model architectures, functioning as synthesized cultural informants, detect, categorize, and quantify latent semantic dissonance between claimed egalitarian ideals and enacted exploitation of others?
\end{itemize}

\section{One-Line Scholarly Contribution}
\begin{quote}
\textbf{The Contribution Formula:} \textit{Before this study}, relationship scholarship assumed that vocal advocacy for equality was an unalloyed indicator of prosocial relationship commitment; \textit{we demonstrate that} parity rhetoric is strategically operationalized as an exploitative lever in the treatment of others that licenses mistreatment; \textit{this matters because} clinical counselors and relationship assessment instruments currently treat progressive fluency as prosocial, blinding them to therapeutic gaslighting and interpersonal predation.
\end{quote}

\section{Pre-Mortem Risk Audit \& Structural Countermeasures}
Conducting a strict \texttt{/premortem} audit reveals three fatal rejection risks and their proactive remedies:
\begin{table}[htbp]
\centering
\small
\begin{tabular}{p{3.5cm}p{5.5cm}p{5.5cm}}
\toprule
\textbf{Potential Fatal Risk} & \textbf{Vulnerability Mechanism} & \textbf{Proactive Hardening Strategy} \\
\midrule
\textbf{1. Conflation of Sincere Equity vs. Weaponized Parity} & Reviewers argue that not all demands for fairness are insincere, risking false-positive classification. & Implement a dual-axis distance metric: isolate linguistic proximity of fairness vocabulary directly preceding documented non-cooperative, asymmetric relationship actions. \\
\textbf{2. LLM Hallucination / Bias as Cultural Informant} & LLMs might project Western socio-linguistic training biases onto naturalistic relationship text. & Triangulate 3 independent model families (OpenAI, Anthropic, Google) with temperature $T=0.1$ and cross-validate against expert human clinical coders. \\
\textbf{3. Ethics of Forum Data Scraping (IRB / PDPA)} & Ethical scrutiny over analyzing public relationship disclosures. & Ensure complete anonymization, cryptographically hash user identifiers (SHA-256), remove all PII, and operate under IRB-exempt public repository guidelines. \\
\bottomrule
\end{tabular}
\caption{Pre-Mortem Vulnerability Analysis and Methodological Safeguards.}
\end{table}

\bibliographystyle{apalike}
\bibliography{library}

\end{document}
